\subsection*{先构造一个容易的破坏过程}
直接生成复杂图像很难，但给图像加噪很容易。扩散方法规定一系列噪声水平，把数据逐渐混入随机噪声；训练时知道原图、噪声和阶段，就能构造监督任务，让模型估计噪声或干净信号。反向生成尝试从噪声出发，逐步执行学习到的更新。

正向公式中的两个系数分别控制信号和噪声强度。它描述的是已知阶段的边际采样，不等于完整算法的所有步骤。模型还需要时间或噪声水平信息，因为同一个观测在不同噪声强度下应有不同解释。

\subsection*{代数可逆与统计可恢复不同}
在标量玩具模型 $x_t=ax_0+b\varepsilon$ 中，若知道真噪声且 $a\ne0$，减掉噪声再除以 $a$ 就能恢复原信号。这里的代数完全确定，正是 Lean 校验的内容。但实际只看到 $x_t$，真噪声未知，多组信号噪声都可能产生同一观测。

若估计噪声误差为 $\delta$，重建信号误差是 $-(b/a)\delta$。当信号系数很小，同样噪声误差会被显著放大。这解释了为什么不能把一次代数反演当成真实生成过程的完整答案。

文本条件影响模型更倾向恢复哪类结构，逐步采样则涉及随机性与具体更新规则。本节建立“容易加噪、学习条件恢复”的思想；训练有效性、生成质量和完整随机过程需要额外分析，不能由一条恢复恒等式推出。
